\documentclass[11pt,a4paper]{article}
\usepackage[margin=24mm]{geometry}
\usepackage[T1]{fontenc}
\usepackage[utf8]{inputenc}
\usepackage{lmodern,microtype}
\usepackage{amsmath,amssymb,amsthm,booktabs,array,xcolor}
\usepackage[hyphens]{url}
\usepackage[colorlinks=true,linkcolor=blue!50!black,urlcolor=blue!50!black,citecolor=blue!50!black]{hyperref}
\hypersetup{pdftitle={A Dense-Complement Route Against the Nemeth--Seberry--Shu Conjecture: Complete Manuscript and Independent Audit},pdfauthor={Showa; Heisei},pdfsubject={Round 10, consolidated original manuscript and independent audit}}
\setlength{\emergencystretch}{3em}
\allowdisplaybreaks
\newcommand{\per}{\operatorname{per}}
\newcommand{\E}{\mathbb E}
\newcommand{\Z}{\mathbb Z}
\newtheorem{theorem}{Theorem}
\newtheorem{lemma}[theorem]{Lemma}
\newtheorem{proposition}[theorem]{Proposition}
\renewcommand{\labelitemi}{$\bullet$}
\title{A Dense-Complement Route Against the\\Nemeth--Seberry--Shu Conjecture\\[7pt]\large Complete Manuscript and Independent Audit}
\author{Showa: main manuscript\qquad Heisei: independent audit\\[3pt]\small Mathematical reasoning tournament, round 10}
\date{Original manuscripts: 16 November 2026\\\small Consolidated edition: 19 November 2026}
\begin{document}
\maketitle
\begin{center}\small
\textbf{Part I. Main manuscript --- Showa}
\end{center}
\begin{abstract}
Let $V(p)$ be the number of distinct permanents of all binary circulants of prime order $p$, with no restriction on their row sums. The 1979 Nemeth--Seberry--Shu conjecture asserts $V(p)=O(p)$. This submission presents a global disproof candidate: $V(p)=\Omega(p^{4/3})$. The argument combines a published uniform dense-complement permanent expansion with an additive-energy contrast, a bounded-change path through subsets, and a prime-order congruence separating row sums. It uses no unproved mathematical hypothesis. Its status is nevertheless \emph{proof candidate submitted for independent audit}, not an accepted resolution or a priority claim. All imported estimates and the limits of the literature audit are identified.
\end{abstract}

\section{Target, scope, and provenance}
Let $P$ be the permutation matrix of translation by one on $\Z_p$. For $S\subseteq\Z_p$, put
\[
 A_S=\sum_{s\in S}P^s,\qquad
 V(p)=\bigl|\{\per A_S:S\subseteq\Z_p\}\bigr|.
\]
The target is the original prime-order assertion
\begin{equation}\label{target}
 V(p)=O(p)\qquad(p\longrightarrow\infty\text{ through primes}).
\end{equation}
Nemeth, Seberry and Shu discuss all $2^p$ binary circulants in their original introduction, not one fixed row sum \cite{NSS}. Cheon--Wanless state this as Conjecture 27 \cite[p.~329]{CW}. The presently inspected status page still labels it open \cite{Lavi}. The constant implicit in $O(p)$ is unrestricted. Consequently, a finite example with more than $p$ values does \emph{not} refute the conjecture.

The separate Resta--Sburlati conjecture concerns a fixed row sum $k$ and predicts $p^{k-2}/k!+O(p^{k-3})$ values \cite{RS,CW}. We do not prove that prediction. Our row sums grow with $p$ and lie close to $p$.

\begin{theorem}[Main disproof candidate]\label{main}
There are constants $c>0$ and $p_0$ such that
\[
 V(p)\ge c p^{4/3}
\]
for every prime $p\ge p_0$. In particular, \eqref{target} is false.
\end{theorem}

\section{The published analytic input}
Write $J_n$ for the all-one matrix. If $C$ is a binary $n\times n$ matrix with all row and column sums $r$, let $Q(C)$ count its all-one $2\times2$ submatrices. These are the four-cycles in its bipartite graph, or the all-zero $2\times2$ submatrices of $J_n-C$.

The following is the Godsil--McKay estimate \cite{GM}, in the shortened form explicitly printed in \cite[p.~321, equation (5)]{CW}. For fixed $\delta>0$ and $r=O(n^{1-\delta})$,
\begin{equation}\label{input}
 \log\per(J_n-C)=B(n,r)+a(n,r)Q(C)+O(r^5/n^5),
\end{equation}
for all such matrices, where
\begin{align}
 a(n,r)&=\frac1{n^4}+\frac{2(2r-1)}{n^5},\label{coef}\\
 B(n,r)&=\log(n!)+n\log(1-r/n)+\frac r{2n}
 +\frac{3r^2-r}{6n^2}+\frac{2r^3-r}{4n^3}\notag\\
 &\hspace{15mm}+\frac{15r^4+70r^3-105r^2+32r}{60n^4}.\label{base}
\end{align}
We use its uniform meaning: on the range $1\le r\le2n^{2/3}$ there is a constant $K$, independent of $n,r,C$, such that the absolute remainder in \eqref{input} is at most $Kr^5/n^5$ for all sufficiently large $n$. This is the fixed choice $\delta=1/3$, with the bound constant $2$. In particular, the estimate is for individual matrices, not an average over random Latin rectangles.

The source PDF's equation, normalization, zero-submatrix definition, and parameter range were visually checked. The degree-only expression $B(n,r)$ cancels in every comparison below. Equation \eqref{input} is \emph{imported}, not a result claimed as new here. It is the only non-elementary input. The logarithm is legitimate because a nonzero regular bipartite graph has a perfect matching; here $n-r>0$.

\section{Two exact combinatorial facts}
For this section $p$ is an odd prime, $S\subseteq\Z_p$, $|S|=r$, and
\[
 C_S=\sum_{s\in S}P^s,\qquad
 c_d(S)=|S\cap(S+d)|,\qquad
 \mathcal E(S)=\sum_{d\in\Z_p}c_d(S)^2.
\]
We use $C_S$ as a sparse complement; the dense matrix of interest is $J_p-C_S=A_{\Z_p\setminus S}$.

\begin{lemma}[Four-cycles and two endpoints]\label{endpoints}
For every $S$,
\begin{equation}\label{Qenergy}
 Q(C_S)=\frac p4\bigl(\mathcal E(S)-2r^2+r\bigr).
\end{equation}
If $2r<p$ and $H=\{0,1,\ldots,r-1\}$, then
\begin{equation}\label{high}
 Q(C_H)=\frac{pr(r-1)(r-2)}6.
\end{equation}
There exists an $r$-subset $L$ such that
\begin{equation}\label{low}
 Q(C_L)\le\frac{pr(r-1)^2(r-2)}{4(p-2)}.
\end{equation}
\end{lemma}
\begin{proof}
For two distinct rows separated by $d$, the number of common ones is $c_d$. Counting unordered row pairs and then column pairs gives
\[
 Q(C_S)=\frac p2\sum_{d\ne0}\binom{c_d}2
       =\frac p4\sum_{d\ne0}(c_d^2-c_d).
\]
Now $c_0=r$ and $\sum_{d\ne0}c_d=r(r-1)$, proving \eqref{Qenergy}. For the interval $H$, the nonzero correlations, apart from $c_0$, are two copies of $1,2,\ldots,r-1$. Hence
\[
 \mathcal E(H)=r^2+2\sum_{j=1}^{r-1}j^2=\frac{2r^3+r}{3},
\]
which yields \eqref{high}.

Choose $S$ uniformly among all $r$-subsets. The energy counts ordered solutions $a-b=c-d$ inside $S$. Among all solutions in $\Z_p$, the numbers having respectively $1,2,3,4$ distinct entries are
\[
 p,\quad 2p(p-1),\quad 2p(p-1),\quad p(p-1)(p-3).
\]
Indeed, the two-entry cases are the two trivial pairings. A three-entry solution has either $a=d$ or $b=c$, giving two arithmetic-progression patterns; these are disjoint for odd $p$. Subtracting these and the one-entry solutions from the total $p^3$ gives the last count. Multiplying by the inclusion probabilities $(r)_j/(p)_j$ and simplifying, for $p\ge5$, gives
\[
 \E\mathcal E(S)=2r^2-r+\frac{r(r-1)^2(r-2)}{p-2}.
\]
The same identity at $p=3$ follows directly. Thus the right-hand side of \eqref{low} is exactly $\E Q(C_S)$. Some $L$ has $Q(C_L)$ no greater than this mean.
\end{proof}

\begin{lemma}[A replacement changes few four-cycles]\label{stepQ}
If $r$-subsets $S,T$ differ by replacing one element, then
\[
 |Q(C_S)-Q(C_T)|\le2pr^2.
\]
\end{lemma}
\begin{proof}
Their bipartite graphs differ in exactly $2p$ edges. Their union is $(r+1)$-regular. A specified edge in that union lies in at most $r^2$ four-cycles: choose one further neighbor of each endpoint, and then check the closing edge. Every four-cycle counted on only one side contains a changed edge. Counting all such possibilities, allowing overcounting, proves the inequality.
\end{proof}

\section{Turning a range and small steps into many values}
\begin{lemma}[Finite-path value count]\label{path}
If real numbers $x_0,\ldots,x_m$ satisfy $|x_m-x_0|\ge h>0$ and $|x_{i+1}-x_i|\le d$ for all $i$, then they take at least $1+h/d$ distinct values.
\end{lemma}
\begin{proof}
Arrange the visited values in increasing order. No gap between successive values can exceed $d$: a part of the original path joining visits on opposite sides would have to cross that gap in one step. Thus the range of $N$ distinct values is at most $(N-1)d$, and is at least $h$.
\end{proof}

Set
\begin{equation}\label{scale}
 R=\lfloor p^{2/3}\rfloor,\qquad R\le r\le2R,
 \qquad F_r(S)=\log\per(J_p-C_S).
\end{equation}
All constants below are independent of the integer $r$ in this range. For sufficiently large $p$, $2r<p$, so Lemma~\ref{endpoints} applies. Its two bounds give
\[
 Q(C_H)-Q(C_L)=\left(\frac16+o(1)\right)pr^3.
\]
Here $Q(C_L)/(pr^3)=O(r/p)$ and the interval's relative error is $O(1/r)$. By \eqref{input}--\eqref{coef},
\begin{align}
 F_r(H)-F_r(L)
 &=a(p,r)\bigl(Q(C_H)-Q(C_L)\bigr)+O(r^5/p^5)\notag\\
 &=\left(\frac16+o(1)\right)\frac{r^3}{p^3}
 \ge\frac{c_0}{p}\label{range}
\end{align}
for some $c_0>0$. The remainder is smaller than the main term by a factor $O(r^2/p^2)=o(1)$, uniformly in \eqref{scale}.

For adjacent $r$-subsets, Lemma~\ref{stepQ} and the same uniform remainder yield
\begin{align}
 |F_r(S)-F_r(T)|
 &\le a(p,r)\,2pr^2+2Kr^5/p^5\notag\\
 &=O(r^2/p^3+r^5/p^5)
 \le Dp^{-5/3}\label{smallsteps}
\end{align}
for a fixed $D>0$. Notice that we did \emph{not} differentiate or cancel the unknown remainder: we bounded the two remainder terms separately.

The graph of $r$-subsets under single-element replacement is connected. Explicitly, pair the elements of $H\setminus L$ with those of $L\setminus H$ and replace them in turn. Applying Lemma~\ref{path} to this path, \eqref{range}, and \eqref{smallsteps}, the logarithms, and hence the permanents, assume at least
\begin{equation}\label{perdegree}
 c_1p^{2/3}
\end{equation}
distinct values among matrices of row sum $p-r$, for some fixed $c_1>0$. Monotonicity along the path is not required.

\section{Separating the row sums}
\begin{lemma}[Prime-order congruence]\label{congruence}
For prime $p$ and any $T\subseteq\Z_p$,
\[
 \per A_T\equiv |T|\pmod p.
\]
\end{lemma}
\begin{proof}
Let the cyclic translation group act by conjugation on the permutations supported by $A_T$. Circulant invariance preserves support. Every orbit has size $1$ or $p$. A fixed permutation commutes with a full $p$-cycle, so it is a translation $x\mapsto x+t$. Exactly $|T|$ translations are supported. The remaining permutations occur in multiples of $p$.
\end{proof}

For distinct $r\in[R,2R]$, the corresponding dense row sums $p-r$ are distinct integers in $\{1,\ldots,p-1\}$. Lemma~\ref{congruence} puts their permanent values in distinct residue classes. Therefore their value sets are disjoint. Summing \eqref{perdegree} over the $R+1$ choices gives
\[
 V(p)\ge(R+1)c_1p^{2/3}\ge c p^{4/3}.
\]
This completes the argument for Theorem~\ref{main}. Since $p^{4/3}/p\to\infty$ along primes, it addresses the original unrestricted $O(p)$ assertion, not merely the stronger and unjustified interpretation $V(p)\le p$.

\section{Computational checks and their limits}
The included C++ program exhaustively computes permanents by subset dynamic programming, four-cycle counts both directly and by \eqref{Qenergy}, the exact mean in Lemma~\ref{endpoints}, and every directed single-element replacement at the following orders:
\begin{center}
\begin{tabular}{rrrr}
\toprule
$p$ & Matrices & $V(p)$ & Directed replacements checked\\
\midrule
3&8&4&12\\
5&32&6&160\\
7&128&9&1,344\\
11&2,048&27&56,320\\
13&8,192&71&319,488\\
\bottomrule
\end{tabular}
\end{center}
All congruences, moment identities, interval formulas in their stated range, and replacement bounds passed. The permanent values are below $13!$, so unsigned 64-bit arithmetic cannot overflow. An independent Python implementation compared subset dynamic programming with Ryser inclusion--exclusion on 25 seeded cases at these orders; every comparison passed.

The same Python file checks energy contrasts and replacement paths at $p=101,1009,10007$, with respective $r=21,100,464$. Their interval four-cycle counts are $134330$, $163155300$, and $165536434448$; the smallest of ten sampled random counts are $38683$, $23628762$, and $11468112063$. These are exact combinatorial checks, \emph{not} computed permanents at large orders. The asymptotic theorem follows from the symbolic argument and the published uniform estimate, not these samples. No effective threshold $p_0$ or numerical value of the imported error constant $K$ is claimed.

\section{G1--G6 audit and submission status}
\begin{description}
\item[G1: Existing name.] The target is the 1979 Nemeth--Seberry--Shu conjecture, conventionally Minc Conjecture 27, not a newly invented assertion.
\item[G2: Original scope.] The original introduction and the later survey both cover all row sums. The proof sums dense row sums precisely because that scope permits it.
\item[G3: Unresolved at selection.] The 2005 survey describes a competing conjecture; the inspected current status page still labels this problem open. Searches on 16 November 2026 located no accepted resolution. This is a bounded status check, not a proof that none exists.
\item[G4: Global and unconditional.] The text supplies a full asymptotic disproof argument, using a published theorem rather than a conjectural expansion. It applies to every sufficiently large prime and needs no experimental assumption. At submission it remains a proof candidate awaiting independent mathematical review; it is not presented as an externally certified solution.
\item[G5: Claimed contribution boundary.] The candidate contribution is the combination of a large four-cycle range, small-step value counting, and congruence-separated row sums. Neither the Godsil--McKay expansion nor affine invariance, elementary moments, or the prime congruence is claimed as a new discovery. Priority for the combination has not been established.
\item[G6: Duplication and falsification checks.] The shared 190-file corpus was searched for the names, distinct-permanent terminology, and Godsil--McKay. Its revision was \texttt{5887cdea4898}. The N\'emeth hits were unrelated N\'emeth--Zim\'anyi references on Pad\'e approximation in files 117/118. Other Godsil and circulant hits concerned unrelated spectral, weighing-matrix, or permanent questions. No matching target or proof was found. The original 1979 paper and the affine-equivalence section of Eades--Praeger--Seberry \cite{EPS} were inspected. The Resta--Sburlati publisher abstract and survey account were read, but not its full article. The Godsil--McKay author scan was partly accessible through indexed text, including its walk-count lemma and extension theorem; direct downloads repeatedly timed out. The actual dependency used here was checked in the downloaded, visually inspected Cheon--Wanless equation. Broader searches did not reveal this dense-row-sum argument, but that does not certify novelty.
\end{description}

The earlier exploratory route using a longer expansion and special walk counts has been retained separately for audit history; it is not needed in this proof. The current argument deliberately avoids that additional dependency. If independent review finds that the uniform form of \eqref{input} or another step fails, Theorem~\ref{main} must be withdrawn rather than re-labelled as proved. No journal acceptance, formal verification, or precedence is asserted.

\clearpage
\begin{thebibliography}{9}
\bibitem{NSS} E. Nemeth, J. Seberry and M. Shu,
\emph{On the distribution of the permanent of cyclic $(0,1)$ matrices},
Utilitas Mathematica \textbf{16} (1979), 171--182.
\href{https://documents.uow.edu.au/~jennie/WEBPDF/071_1979.pdf}{Author-hosted original scan}.
\bibitem{CW} G.-S. Cheon and I. M. Wanless,
\emph{An update on Minc's survey of open problems involving permanents},
Linear Algebra and its Applications \textbf{403} (2005), 314--342.
\href{https://doi.org/10.1016/j.laa.2005.02.030}{DOI};
\href{https://users.monash.edu.au/~iwanless/papers/permsurveyLAA.pdf}{Author PDF}.
\bibitem{GM} C. D. Godsil and B. D. McKay,
\emph{Asymptotic enumeration of Latin rectangles},
Journal of Combinatorial Theory, Series B \textbf{48} (1990), 19--44.
\href{https://doi.org/10.1016/0095-8956(90)90128-M}{DOI};
\href{https://users.cecs.anu.edu.au/~bdm/papers/LatinRectangles.pdf}{Author-hosted corrected scan}.
\bibitem{EPS} P. Eades, C. E. Praeger and J. R. Seberry,
\emph{Some remarks on the permanents of circulant $(0,1)$ matrices},
Utilitas Mathematica \textbf{23} (1983), 145--159.
\href{https://documents.uow.edu.au/~jennie/WEBPDF/092_1983.pdf}{Author-hosted original scan}.
\bibitem{RS} G. Resta and G. Sburlati,
\emph{On the number of different permanents of some sparse $(0,1)$-circulant matrices},
Linear Algebra and its Applications \textbf{375} (2003), 197--209.
\href{https://doi.org/10.1016/S0024-3795(03)00649-9}{DOI and abstract}.
\bibitem{Lavi} Y. Lavi,
\emph{Few permanents of prime circulants}, Minc's list, Conjecture 27.
\href{https://yair-lavi.com/mincs-list/conjecture-27/index.html}{Status and prior-progress page}, inspected 16 November 2026.
\end{thebibliography}

\clearpage
\phantomsection
\addcontentsline{toc}{part}{Part II. Independent audit --- Heisei}
\begin{center}
{\Large\bfseries Part II. Independent Audit}\par\medskip
{\large Independent Audit of Showa's Minc--27 Candidate v2}\par\smallskip
Heisei\par
16 November 2026
\end{center}
\medskip
\noindent\textit{Editorial note.} The following is the complete independent audit
from the original archive. Its cross-references and bibliography keys have been
renamed for this single-file edition; its mathematical content and qualification
of the verdict are preserved. References to auxiliary scripts and the reviewed
snapshot describe the original evidence package and are not compilation inputs.
\medskip
\setcounter{section}{0}
\setcounter{equation}{0}
\renewcommand{\thesection}{A.\arabic{section}}
\renewcommand{\theequation}{A.\arabic{equation}}
\renewcommand{\theHsection}{audit.\arabic{section}}
\renewcommand{\theHequation}{audit.\arabic{equation}}
\renewcommand{\refname}{References for the Independent Audit}
\begin{abstract}
This is an adversarial audit of a fixed snapshot, not a new research submission or an assertion of priority. No error was found in the v2 derivation of $V(p)=\Omega(p^{4/3})$ from the published dense-complement permanent expansion. The relevant source allows growing row deficit, and the original Godsil--McKay paper explicitly makes its section-6 error constants independent of that deficit. We give exact elementary derivations, a uniform threshold calculation, and fresh exhaustive finite diagnostics. These findings address mathematical completeness of the candidate, not G5 originality, exhaustive literature priority, formal proof certification, or independent rederivation of the imported analytic theorem.
\end{abstract}

\section{Reviewed object and precise target}
The reviewed text is \texttt{proof\_candidate\_reviewed.txt}, copied without changing Showa's original. Its SHA-256 is
\begin{center}\small\ttfamily
c7016e89d476c96a7e9da5926cac16df4109\\f4c1dd6c71120f8dffdae0fb0e0b.
\end{center}
It is explicitly the simplified v2 argument. Claims or flaws concerning the earlier epsilon-walk route do not automatically apply to this snapshot.

Let $P$ be a cyclic shift of order $p$, and put $C_S=\sum_{s\in S}P^s$ for $S\subseteq\mathbb Z_p$. The target is
\[
 V(p)=\#\{\per C_S:S\subseteq\mathbb Z_p\}=O(p)
 \qquad(p\longrightarrow\infty\text{ through primes}).
\]
Nemeth--Seberry--Shu's original opening discussion considers all binary circulants, and Cheon--Wanless state this as Conjecture 27 \cite{audit:NSS,audit:CW}. The target is neither the inequality $V(p)\le p$ nor a fixed-row-sum conjecture. The separate rival Conjecture $27'$ in \cite{audit:CW} fixes the row sum and is not settled by v2.

\section{The imported theorem and uniformity}
The actual image of \cite[p.~321, Eq.~(5)]{audit:CW} was inspected, not merely an OCR rendering. With $C$ a binary $n\times n$ matrix having every row and column sum $r$, its content is
\begin{equation}\label{audit:expansion}
 \log\per(J-C)=B(n,r)+a(n,r)Q(C)+E_C,
 \qquad |E_C|\le K\frac{r^5}{n^5},
\end{equation}
where $Q(C)$ counts unordered choices of two rows and two columns on which all four entries are one, and
\[
 a(n,r)=\frac1{n^4}+\frac{2(2r-1)}{n^5}.
\]
The common scalar $B$ is the log factorial and logarithmic prefactor plus the terms depending only on $n,r$ through denominator $n^4$. Its exact value cancels from every comparison in v2. The displayed error denominator really is $n^5$, and the structural statistic is a count of zero $2\times2$ submatrices of $J-C$, equivalently $Q(C)$.

The source range is $r=O(n^{1-\delta})$, with fixed $\delta>0$, for all regular matrices in question. Fixing $\delta=1/3$ and the growth bound $r\le2n^{2/3}$ gives constants $K\ge1$ and $n_0$ such that \eqref{audit:expansion} holds for every admissible $n\ge n_0$, $r$, and $C$. Neither $K$ nor $n_0$ is allowed to vary with the selected matrix or with $r$ inside this fixed regime. They may depend on the fixed growth regime. Calling $K$ ``absolute'' should be understood this way; a stronger assertion uniform over all $\delta$ is unnecessary.

\paragraph{Original-source cross-check.}
The web PDF extraction of Godsil--McKay \cite[\S6, pp.~36--39]{audit:GM} was also read. At the start of that section the authors state: ``the implied constant is independent of $k$.'' The section fixes $k=O(n^{1-\delta})$. Theorems 6.2 and 6.3 give the individual-rectangle expansion, not just an average over random rectangles. Section 2 identifies extensions with complement perfect matchings. Every regular bipartite graph has a one-factorization, by iterated Hall matching, so that formulation covers every binary regular $C$ used here. The average formula in their Theorem 6.4 is not needed.

Theorem 6.3 has the same $s/n^4+2(2k-1)s/n^5$ structural terms, with additional sixth-denominator terms and error. For a degree-$k$ bipartite graph on $2n$ vertices, the elementary bounds $s=O(nk^3)$, $h=O(nk^5)$, and $b=O(nk^4)$ for squares, hexagons, and copies of $K_{2,3}$ make the dropped structural terms $O(k^5/n^5)$. The dropped scalar fifth-degree term has the same order; later terms are smaller for $k/n\to0$. This corroborates the shortened formula in \cite{audit:CW}.

The original PDF's local download timed out. Its web extraction, including the uniformity statement and theorem statements, was accessible. The full symbolic integration leading to \cite[Theorem 6.2]{audit:GM} was not independently re-executed. This is ordinary reliance on a published theorem, not a newly assumed conjectural estimate.

\section{Independent exact counting audit}
For $|S|=r$, set $c_d=|S\cap(S+d)|$ and $\mathcal E(S)=\sum_dc_d^2$. Direct row-pair counting gives
\begin{align}
 Q(C_S)&=\frac p2\sum_{d\ne0}\binom{c_d}{2}
       =\frac p4\bigl(\mathcal E(S)-2r^2+r\bigr).\label{audit:Q}
\end{align}
Indeed, $\sum_dc_d=r^2$ and $c_0=r$. The factor $p/2$ removes the double counting of ordered row pairs, and the binomial counts unordered column pairs. Thus $Q$ is neither an oriented cycle count nor a count multiplied by eight.

For $H=\{0,\ldots,r-1\}$ and $2r<p$, the nonzero difference multiplicities are $r-j$ at both $j$ and $-j$, $1\le j<r$. Consequently
\[
 \mathcal E(H)=r^2+2\sum_{j=1}^{r-1}j^2=\frac{2r^3+r}{3},
 \qquad Q(C_H)=\frac{pr(r-1)(r-2)}6.
\]
The condition $2r<p$ prevents the positive and negative difference supports from overlapping.

For uniform $r$-subsets of an odd prime group, classify ordered solutions of $a-b=c-d$ by their number $j$ of distinct entries. The counts are
\[
 (N_1,N_2,N_3,N_4)
 =\bigl(p,\,2p(p-1),\,2p(p-1),\,p(p-1)(p-3)\bigr).
\]
For $j=2$ the two possibilities pair $a=b,c=d$ or $a=c,b=d$. For $j=3$ the repeated entry is $a=d$ or $b=c$, giving two arithmetic-progression patterns. Oddness ensures that each has the stated count; their common all-equal case is removed. Subtracting these counts from the $p^3$ total solutions gives $N_4$. Multiplying $N_j$ by $(r)_j/(p)_j$ and simplifying yields
\begin{align}
 \E\mathcal E(S)&=2r^2-r+\frac{r(r-1)^2(r-2)}{p-2},\\
 \E Q(C_S)&=\frac{pr(r-1)^2(r-2)}{4(p-2)}.\label{audit:average}
\end{align}
There exists an $r$-subset $L$ no larger than this latter average. No randomness or concentration property is subsequently required of the chosen $L$.

If $S,T$ differ by one replacement, their incidence graphs differ in exactly $2p$ edges and their union is $(r+1)$-regular. For each changed edge, choosing another neighbor of each endpoint gives at most $r^2$ four-cycles of that union containing it. Every cycle present on only one side contains a changed edge. Overcounting is harmless, giving the uniform bound
\begin{equation}\label{audit:swap}
 |Q(C_S)-Q(C_T)|\le2pr^2.
\end{equation}
This includes both deleted and inserted cycles. It does not silently presume that swaps change only one edge of a circulant.

\section{A quantified all-large-primes calculation}
Let $R=\lfloor p^{2/3}\rfloor$ and $R\le r\le2R$. Increase $K$ in \eqref{audit:expansion} to at least one if necessary. Take every prime $p$ satisfying
\begin{equation}\label{audit:threshold}
 p\ge n_0,\qquad p\ge512,\qquad p^{2/3}\ge768K.
\end{equation}
Then $4\le r\le p/4$, in particular $2r<p$. From the exact endpoint and average formulas,
\begin{align*}
 Q(C_H)-Q(C_L)
 &\ge Q(C_H)\left(1-\frac{3(r-1)}{2(p-2)}\right)\\
 &\ge\frac12Q(C_H)\ge\frac{pr^3}{48}.
\end{align*}
Here $(r-1)(r-2)\ge r^2/4$ and $3(r-1)\le p-2$ suffice. The logarithmic endpoint separation $h$ therefore obeys
\begin{align}
 h&:=\log\per(J-C_H)-\log\per(J-C_L)\nonumber\\
 &\ge\frac{r^3}{48p^3}-2K\frac{r^5}{p^5}
 \ge\frac{r^3}{96p^3}\ge\frac1{768p}.\label{audit:gap}
\end{align}
The middle inequality uses $2Kr^2/p^2\le8Kp^{-2/3}\le1/96$. The last uses $R\ge p^{2/3}/2$. This explicitly handles the same $n_0$ and the same $K$ for all chosen row deficits. The main endpoint signal is order $p^{-1}$; the two endpoint errors together are order $p^{-5/3}$, not comparable to or larger than that signal.

Also $a(p,r)\le2p^{-4}$ for $r\le p/4$. For a swap, \eqref{audit:expansion} and \eqref{audit:swap} give
\begin{align}
 |\Delta\log\per|
 &\le\frac{4r^2}{p^3}+\frac{2Kr^5}{p^5}
 \le (16+64K)p^{-5/3}=:d.\label{audit:steps}
\end{align}
Even though the remainder need not itself vary smoothly under a swap, the difference of two remainders is bounded by twice their individual bound. That is all the proof uses.

\section{Distinct values, paths, and residues}
The graph of $r$-subsets with one-element replacements is connected: replace elements of $H\setminus L$ by elements of $L\setminus H$ one at a time. For any finite real sequence having step size at most $d>0$, let its $N$ distinct values in increasing order be $v_1<\cdots<v_N$. Every consecutive gap is at most $d$, since otherwise the sequence could not cross that gap to visit both sides. Therefore the endpoint difference is at most $(N-1)d$, whether or not the sequence is monotone. Applying \eqref{audit:gap}--\eqref{audit:steps} yields at least
\[
 1+\frac{p^{2/3}}{768(16+64K)}
\]
distinct logarithms, and hence distinct permanents, at each dense row sum $p-r$. Those permanents are positive because a regular bipartite graph of positive degree has a perfect matching.

For a binary circulant $A_T$ of prime order $p$, simultaneous cyclic conjugation acts on the supported permutations. Its orbits have sizes one or $p$. The fixed permutations commute with a full $p$-cycle and therefore are translations. Exactly $|T|$ translations are supported. It follows that
\[
 \per A_T\equiv |T|\pmod p.
\]
The integers $p-r$ for $R\le r\le2R<p$ are distinct modulo $p$, so no permanent value can be counted at two different selected row sums. Summing over $R+1\ge p^{2/3}$ values of $r$ gives
\begin{equation}\label{audit:bound}
 \boxed{V(p)\ge\frac{p^{4/3}}{768(16+64K)}}
\end{equation}
for every prime satisfying \eqref{audit:threshold}. Thus the two $p^{2/3}$ factors really multiply. The argument is not just a lower bound on path length, and it does not count repeated permanent values as new values. The unknown but fixed $K,n_0$ affect the positive constant and starting threshold, not the exponent or the contradiction to $O(p)$.

\section{Fresh finite checks and limits of the verdict}
The independent \texttt{audit\_checks.py} imports no Showa code or data. Integer dynamic programming computes every permanent at $p=3,5,7,11$. Direct unordered row-pair and column-pair counting is compared with \eqref{audit:Q}; exact rational averaging is compared with \eqref{audit:average}. The script also enumerates the ordered quadruple patterns, every directed one-element swap, all nonwrapping interval endpoints, and the prime congruence. Results are saved in \texttt{audit\_checks.json} and \texttt{audit\_checks.log}.
\begin{center}
\begin{tabular}{rrrr}
\toprule
$p$ & All subsets & Directed swaps & Distinct permanents\\\midrule
3&8&12&4\\
5&32&160&6\\
7&128&1344&9\\
11&2048&56320&27\\\bottomrule
\end{tabular}
\end{center}
All checks passed. A further 21,840 finite real-sequence tests of the path inequality passed, including nonmonotone and repeated-value sequences. These finite tests can expose normalization and implementation mistakes but do not establish a uniform asymptotic error bound or determine $n_0$. No inference from a small-prime count to an asymptotic contradiction is used.

\paragraph{Audit conclusion.}
At the level of the mathematical derivation, v2 is complete relative to the cited published expansion: \eqref{audit:bound} contradicts the exact original all-row-sums conjecture. No missing elementary step or mismatch in the imported theorem's range was found. This is stronger than saying that a finite computation exceeded $p$, but it is not external acceptance or formal certification. This auditor did not reconstruct all analytic calculations in the cited theorem, nor prove that the combination is novel. G5 originality and a sufficiently broad G6 priority review must be evaluated separately. A current research page still labels the original conjecture open \cite{audit:Lavi}; that status is contextual evidence, not a substitute for proof review or a guarantee of priority.

\begin{thebibliography}{9}
\bibitem{audit:NSS} E. Nemeth, J. Seberry and M. Shu, \emph{On the distribution of the permanent of cyclic $(0,1)$ matrices}, Utilitas Mathematica 16 (1979), 171--182.\newline
\url{https://documents.uow.edu.au/~jennie/WEBPDF/071_1979.pdf}.
\bibitem{audit:CW} G.-S. Cheon and I. M. Wanless, \emph{An update on Minc's survey of open problems involving permanents}, Linear Algebra and its Applications 403 (2005), 314--342, especially pp.~321, 329--330. DOI: \href{https://doi.org/10.1016/j.laa.2005.02.030}{10.1016/j.laa.2005.02.030}.\newline
\url{https://users.monash.edu.au/~iwanless/papers/permsurveyLAA.pdf}.
\bibitem{audit:GM} C. D. Godsil and B. D. McKay, \emph{Asymptotic enumeration of Latin rectangles}, Journal of Combinatorial Theory, Series B 48 (1990), 19--44. DOI: \href{https://doi.org/10.1016/0095-8956(90)90128-M}{10.1016/0095-8956(90)90128-M}.\newline
\url{https://users.cecs.anu.edu.au/~bdm/papers/LatinRectangles.pdf}.
\bibitem{audit:Lavi} Y. Lavi, \emph{Minc's list: Conjecture 27}, research problem page, accessed 16 November 2026.\newline
\url{https://yair-lavi.com/mincs-list/conjecture-27/index.html}.
\end{thebibliography}
\end{document}
